\documentclass[12pt, twoside]{article}
\input{../preamble}

\fancyhead[LO]{La Scuola d'Italia / Dr.\ Huson / IB Math \\* 8 May 2026}
\fancyhead[RO]{\fbox{\begin{minipage}{6cm}\footnotesize First \& Last name:\\[0.5cm]\end{minipage}}}
\fancyfoot[C]{\thepage}

\begin{document}
\subsubsection*{7.9 PreQuiz: Exponential and Logarithmic Word Problems
(calculator)}

\begin{enumerate}[itemsep=0.15cm]

%--- Problem 1: Population growth (continuous) ---
\item[1.] A bacteria culture begins with $200$ cells. The number of
cells $N$, $t$ hours after the start of the experiment, is modelled by
$$N(t) \;=\; 200\, e^{\,k t}, \qquad k > 0.$$
After $2$ hours, the culture contains $800$ cells.

\begin{enumerate}[itemsep=0.15cm]
  \item Write down the value of $N$ when $t = 0$.

  \item Show that $k = \dfrac{1}{2}\ln 4$.

  \item Find the number of cells after $5$ hours, to the nearest
  integer.

  \item Find the time, in hours, at which the number of cells first
  reaches $5000$. Give your answer correct to three significant
  figures.

  \item Find the doubling time of the culture. Give your answer correct
  to three significant figures.
\end{enumerate}

\begin{tikzpicture}
  \draw (-0.6,0) rectangle (11,11);
  \foreach \y in {1,2,...,10} \draw[dotted] (0,\y)--(10.5,\y);
\end{tikzpicture}

\newpage

%--- Problem 2: Compound interest (annual + continuous comparison) ---
\item[2.] Marco deposits \$3000 in a savings account that pays interest
of $4.5\,\%$ per annum, compounded annually. No further deposits or
withdrawals are made.

\begin{enumerate}[itemsep=0.15cm]
  \item Show that the value of the account at the end of the first
  year is \$3135.

  \item Find the value of the account at the end of $6$ years, correct
  to the nearest dollar.

  \item Find the smallest number of full years required for the value
  of the account to first exceed \$5000.

  \item Suppose that, instead, the interest is compounded continuously
  at a nominal annual rate of $4.5\,\%$, so that the value after $t$
  years is $A(t) = 3000\, e^{\,0.045\, t}$. Find the value of this
  account at the end of $6$ years, correct to the nearest cent, and
  state the difference between the two accounts at that time.
\end{enumerate}

\begin{tikzpicture}
  \draw (-0.6,0) rectangle (11,10);
  \foreach \y in {1,2,...,9} \draw[dotted] (0,\y)--(10.5,\y);
\end{tikzpicture}

\newpage

%--- Problem 3: Radioactive decay (cesium-137) ---
\item[3.] Cesium-137 is a radioactive isotope used in some medical
treatments. The mass remaining, in grams, $t$ years after a sample is
prepared is modelled by
$$M(t) \;=\; M_{0}\, \!\left(\dfrac{1}{2}\right)^{\!t/T},$$
where $M_{0}$ is the initial mass and $T$ is the half-life. For
cesium-137, $T = 30$ years. A sample is prepared with $M_{0} = 64$
grams.

\begin{enumerate}[itemsep=0.15cm]
  \item Find the mass remaining after $50$ years. Give your answer
  correct to three significant figures.

  \item Find the time at which the mass has decayed to $8$ grams.

  \item Find, to the nearest year, the time after which only $5\,\%$
  of the original mass remains.
\end{enumerate}

\begin{tikzpicture}
  \draw (-0.6,0) rectangle (11,12);
  \foreach \y in {1,2,...,11} \draw[dotted] (0,\y)--(10.5,\y);
\end{tikzpicture}

\newpage

%--- Problem 4: Newton's cooling (coffee) ---
\item[4.] A cup of coffee is poured. Its temperature $T$, in degrees
Celsius, $t$ minutes after pouring, is modelled by
$$T(t) \;=\; 22 \;+\; 64\, e^{-k t}, \qquad k > 0.$$
After $4$ minutes, the temperature of the coffee is $54\,^{\circ}$C.

\begin{enumerate}[itemsep=0.15cm]
  \item State the temperature of the coffee at the moment it is
  poured.

  \item State the temperature of the room.

  \item Show that $k = \dfrac{1}{4}\ln 2$.

  \item Find the temperature of the coffee after $10$ minutes. Give
  your answer correct to three significant figures.

  \item The coffee is comfortable to drink when its temperature
  reaches $30\,^{\circ}$C. Find the time at which this occurs.
\end{enumerate}

\begin{tikzpicture}
  \draw (-0.6,0) rectangle (11,10);
  \foreach \y in {1,2,...,9} \draw[dotted] (0,\y)--(10.5,\y);
\end{tikzpicture}

\end{enumerate}
\end{document}
